\chapter{Conditional Hypotheses for F1--F3 Arithmetic}
\label{v4-arith:w3-comparison-A}

Several F1--F3 arithmetic reductions are correct as equivalences, but
some component subconditions require open inputs:
the Fargues--Scholze-to-Emerton--Gee categorical equivalence
(F1.a, an unproved comparison hypothesis), the arithmetic
Arinkin--Gaitsgory conjecture
(\(\mathrm{L}_{\mathrm{ACD}}.1\)--\(\mathrm{L}_{\mathrm{ACD}}.3\)),
and the formal-Noetherian extension of Gaitsgory--Rozenblyum
(R1.formal). The exact logical form is the conditional system
determined by
Section~\ref{v4-arith:app:C:residual-table-w2},
Chapter~\ref{v4-arith:platonic-rh-completion}, and
Chapter~\ref{v4-arith:rh-reformulation}.

\section{Conditional Statements}
\label{v4-arith:w3-comparison-A:list}

For each reduction the displayed hypothesis is part of the theorem, not
an afterthought to the proof.

\subsection{F1--F3 in Chapter~\ref{v4-arith:closure}}
\label{v4-arith:w3-comparison-A:ch22}

\paragraph{F1.a domain: Fargues--Scholze local geometrisation versus categorical equivalence.}
Theorem~\ref{v4-arith:thm:F1-reduction} requires
an unproved comparison hypothesis for F1.a, the local categorical
equivalence
\[
\mathrm{D}^{\mathrm{sph\text{-}fin}}([GL_{2}(\mathbb{Q}_{p})])
\simeq
\mathrm{IndCoh}_{\mathrm{Nilp}}(\mathcal{X}^{\mathrm{EG},p}_{GL_{2}})
\]
(where \(\mathcal{X}^{\mathrm{EG},p}_{GL_2}\) is the
Emerton--Gee moduli stack at the prime \(p\),
\(\mathrm{IndCoh}_{\mathrm{Nilp}}\) denotes ind-coherent sheaves
with nilpotent singular support in the sense of Arinkin--Gaitsgory,
and \(\mathrm{D}^{\mathrm{sph\text{-}fin}}\) the spherical-finite
derived category of smooth \(GL_2(\mathbb{Q}_p)\)-representations).
Fargues--Scholze arXiv:2102.13459 supplies the local geometrisation,
spectral action, and Langlands-parameter structure for \(GL_2\); the
displayed categorical equivalence is not contained there.
Dotto--Emerton--Gee arXiv:2603.26887 gives a fully-faithful-only,
fixed-central-character functor for \(GL_{2}(\mathbb{Q}_{p})\) at
\(p\geq 5\). The authoritative reference for the current F1.a form
is Section~\ref{v4-arith:app:C:residual-table-w2} and
Theorem~\ref{v4-arith:thm:F1-reduction}.

\paragraph{Emerton--Gee deformation stack.}
Emerton--Gee arXiv:1908.07185 constructs the deformation-theoretic setting
(\emph{Moduli stacks of \'etale \((\varphi,\Gamma)\)-modules and the
existence of crystalline lifts}, Annals of Mathematics Studies).
See App.~A \S\ref{v4-arith:app:A:langlands}.

\paragraph{Coherent Springer theory.}
Ben-Zvi--Chen--Helm--Nadler arXiv:2010.02321 treats coherent Springer theory (\emph{Coherent Springer theory and the categorical
Deligne--Langlands correspondence}, Invent.~Math.~235 (2024)).
See App.~A \S\ref{v4-arith:app:A:BZN-coherent-Springer}.

\subsection{Adelic Categorical Descent}
\label{v4-arith:w3-comparison-A:ch27}

\paragraph{Function-field transport and \((L_2)\).}
Theorem~\ref{v4-arith:acd:thm:function-field-transport} requires
an unproved comparison hypothesis. The assertion \((L_{1})\) follows
from Raskin's \(\infty\)-categorical construction independently of
function-field Langlands; the assertion \((L_{2})\) requires the
local categorical comparison F1.a (an unproved comparison hypothesis
via Dotto--Emerton--Gee) and Bernstein-centre diamond gluing.

\paragraph{\((L_2)\) reduction: Fargues--Scholze domain.}
Theorem~\ref{v4-arith:acd:thm:L2-reduction} requires
an unproved comparison hypothesis. The comparison functors
\(\Phi^{\mathrm{loc}}_{v}\) identifying \(\mathcal{C}_{v}\) with
\(\mathrm{IndCoh}_{\mathrm{Nilp}}(\mathcal{X}^{\mathrm{EG},v}_{GL_{2}})\)
are not supplied by Fargues--Scholze arXiv:2102.13459; that source
gives the local geometrisation and spectral action. The local
categorical comparison is the residual open input.

\paragraph{F1 reduction: two open inputs.}
Theorem~\ref{v4-arith:acd:thm:F1-closure} conditionally states that F1 reduces to F1.a
plus (BC-diamond-glue). Since F1.a remains
an unproved comparison hypothesis, the reduction leaves F1 conditional
on two independent open inputs: F1.a and (BC-diamond-glue).

\subsection{Arithmetic Ben-Zvi--Nadler}
\label{v4-arith:w3-comparison-A:ch28}

\paragraph{F3.BZN\({}^{\mathrm{Ar}}\): procyclic shadow versus full Weil parameter stack.}
Conjecture~\ref{v4-arith:bzn:thm:main} is conjectural.
The comparison target for F3.BZN\({}^{\mathrm{Ar}}\) is the full
Weil/\((\varphi,\Gamma)\) parameter stack
\(\mathrm{LocSys}^{W,\varphi\Gamma}_{GL_{2}}\), not the procyclic
shadow
\(\mathrm{Map}(B\widehat{\mathbb{Z}},\mathcal{X}^{\mathrm{Sel}}_{\bar\rho,\det,\mathcal{L}})\),
which forgets Frobenius and unramified data. The theorem is
conditional on F1 and F2. Proposition~\ref{v4-arith:bzn:prop:EG-to-LocSys}
carries an unproved comparison hypothesis; the closed-immersion
property of Proposition~\ref{v4-arith:bzn:prop:EG-image-flat} is
conjectural, since the closed-immersion property is
unavailable on the inertia-only shadow.

\subsection{Structural Reduction Theorems}
\label{v4-arith:w3-comparison-A:ch30}

\paragraph{\(\mathrm{L}_{\mathrm{ACD}}\) equivalence and subcondition forms.}
Theorem~\ref{v4-arith:sw:thm:LACD-split} is proved under its stated hypotheses: the
adelic categorical descent condition \(\mathrm{L}_{\mathrm{ACD}}\) is
equivalent to the conjunction of subconditions
\(\mathrm{L}_{\mathrm{ACD}}.1\) through \(\mathrm{L}_{\mathrm{ACD}}.4\).
The equivalence itself is a logical factorisation. The subconditions
\(\mathrm{L}_{\mathrm{ACD}}.1\)--\(\mathrm{L}_{\mathrm{ACD}}.3\) remain
conjectural; \(\mathrm{L}_{\mathrm{ACD}}.4\) is the
combination of Kneser's strong approximation theorem for \(SL_{2}\)
and Tate's determinant idele-class quotient for \(GL_{2}\). The
residual content is recorded in the
subcondition forms.

\paragraph{F3.BZN\({}^{\mathrm{Ar}}\) equivalence via Gaitsgory--Rozenblyum.}
Theorem~\ref{v4-arith:sw:thm:BZN-Ar-split} is proved under its stated hypotheses:
F3.BZN\({}^{\mathrm{Ar}}\), assembled through Gaitsgory--Rozenblyum
DAG Vol.~II, is equivalent to the conjunction of R1, R2, and R3.
The equivalence is a valid logical factorisation. R1 holds as a
theorem for smooth finite-type stacks (Gaitsgory--Rozenblyum) and
is conjectural for the formal-Noetherian Emerton--Gee
setting (R1.formal); R2 is an unproved comparison hypothesis; R3 is
definitional. The R1 fine structure and the honest residual are
recorded in Chapter~\ref{v4-arith:bzn-sub} and
Corollary~\ref{v4-arith:bzn-sub:cor:F3-BZN-closure-path}.

\paragraph{Arithmetic Positselski C1--C5.}
Chapter~\ref{v4-arith:rh-reformulation} treats the arithmetic Positselski
equivalence C1--C5 as five subconditions, of which C1--C4 are
conditional on explicit hypotheses and C5 (functoriality) is the
named residual. The arithmetic Heisenberg equivalence is restricted to
the abelian-Heisenberg sector; C5 is the functoriality condition for
the full claim.

\subsection{\(\mathrm{L}_{\mathrm{ACD}}\) Subconditions}
\label{v4-arith:w3-comparison-A:ch31}

\paragraph{\(\mathrm{L}_{\mathrm{ACD}}.0\): archimedean stratum construction.}
Theorem~\ref{v4-arith:lacd:thm:L0-closed} is proved under its stated hypotheses: the
archimedean local category \(\mathcal{C}_{\infty}\), distinguished
unit, and Langlands--Vogan classification are constructed
from the theorems of Harish-Chandra (1966), Vogan (1981), and
Wallach (1988). Thus \(\mathrm{L}_{\mathrm{ACD}}.0\) is independent of
the global \(\mathrm{L}_{\mathrm{ACD}}\) comparison.

\paragraph{\(\mathrm{L}_{\mathrm{ACD}}.1\): the F2.c hypothesis.}
Theorem~\ref{v4-arith:lacd:thm:L1-closed} is conditional
on F2.c: the assertion \(\mathrm{L}_{\mathrm{ACD}}.1\)
uses the restricted tensor theorem
Theorem~\ref{v4-arith:tate:thm:F2c-main} on the central idempotent
pointed-unit \(E_{2}\) domain. Hence
Theorem~\ref{v4-arith:lacd:thm:L1-closed} is conditional, with F2.c as its displayed
hypothesis.

\paragraph{\(\mathrm{L}_{\mathrm{ACD}}.2\): local Hecke commutativity, not global functoriality.}
Theorem~\ref{v4-arith:lacd:thm:L2-closed} is proved under its stated hypotheses: formal
tensor-factor commutativity of local Hecke endofunctors, once the
restricted tensor product is constructed. The domain remark
Remark~\ref{v4-arith:lacd:rem:L2-formal-domain} records that global Hecke
functoriality is not part of the statement; the
\(\mathrm{L}_{\mathrm{ACD}}.2\) label refers to the local
commutativity subcondition only.

\paragraph{\(\mathrm{L}_{\mathrm{ACD}}\): Bernstein-centre compatibility.}
Corollary~\ref{v4-arith:lacd:cor:LACD-residual} assumes F2.c, the
\(\infty\)-categorical Bernstein lift, the finite-prime local
categorical comparison, the archimedean boundary datum, and the
global identification with the automorphic category.
Under these hypotheses, the remaining Bernstein-centre condition
is (BC-diamond-glue). Its local construction, cross-prime gluing,
and global action compatibility are the three requirements in
Theorem~\ref{v4-arith:bcglue-k5:thm:decomposition}.
The local categorical comparison F1.a remains an input.
Bernstein-centre gluing alone does not supply the other hypotheses.

\subsection{BZN Formal Noetherian and EG Derived}
\label{v4-arith:w3-comparison-A:ch32}

\paragraph{R2: the Selmer-derived target and its formal geometry.}
Fix a residual representation \(\bar\rho\), determinant, finite
ramification set, and local deformation conditions \(\mathcal L\).
\Cref{v4-arith:bzn-sub:thm:R2-main} defines the Selmer-derived
stack \(\mathcal X^{\mathrm{Sel}}_{\bar\rho,\det,\mathcal L}\)
as a derived fibre product of global and local deformation functors.
Linearization gives its Selmer tangent complex.
The local conditions must be derived formal algebraic substacks
of the full Weil/\((\varphi,\Gamma)\) parameter stacks.
Formal Noetherianness, quasi-smoothness, and compatibility of
nilpotent singular support with the Selmer cotangent complex
remain hypotheses for the chosen deformation problem.
A product or colimit presentation alone does not establish them.

\paragraph{R1: extension of the finite-type trace theorem.}
The Gaitsgory--Rozenblyum trace theorem for smooth finite-type
stacks is the finite-type input R1.fin-type.
Its extension to the Selmer-derived target requires R1.formal,
including formal-loop continuity and the flat-locus comparison.
\Cref{v4-arith:bzn-sub:thm:R1-formal-reduction} states this
conditional reduction. The formal-Noetherian hypotheses on the
arithmetic target remain separate from the finite-type theorem.

\subsection{Loop Convention and Wild Primes}
\label{v4-arith:w3-comparison-A:ch33}

\paragraph{R3: comparison of arithmetic loop spaces.}
Let \(\mathcal X\) be a formal Noetherian derived algebraic stack
with cotangent complex of finite Tor-amplitude and a
\(t\)-bounded Postnikov tower.
\Cref{v4-arith:bzn-rcompletion:thm:R3-three-agree} compares its
three loop constructions under three additional assumptions:
\(B\widehat{\mathbb Z}\)-mapping stacks are computed by the
profinite limit of finite cyclic mapping stacks, formal completion
commutes with the relevant derived tensor products, and the perfect
locus generates the same loop object after ind-completion.
These assumptions give the comparison equivalences.
They do not identify the inertia shadow with the full
Weil/\((\varphi,\Gamma)\) parameter stack.

\paragraph{R4: large-prime wild Hochschild vanishing.}
For \(p\ge5\),
Corollary~\ref{v4-arith:bzn-rcompletion:thm:R4-correction-vanishes-p-large}
gives positive-degree vanishing only under
Hypothesis~\ref{v4-arith:bzn-rcompletion:thm:bez-ss-large-p},
which specifies an actual local arithmetic trace equivalent to a
module in degree zero. Semisimplicity of an unspecified category
or a numerical prime threshold does not supply this comparison.
At two and three, the arithmetic local models and Ext-to-trace maps
remain open; their full assembly must respect the four-term tensor
formula of Lemma~\ref{v4-arith:wild:four-term-tensor}.

\subsection{BC-Diamond-Glue Decomposition}
\label{v4-arith:w3-comparison-A:ch35}

\paragraph{Local construction, gluing, and global action.}
The construction data in
\Cref{v4-arith:bcglue-k5:thm:decomposition} give
\[
\text{(BC-diamond-glue)}
\;\Leftrightarrow\;
\text{(BCglue.local)}\wedge\text{(BCglue.glue)}
\wedge\text{(BC-diamond-glue.core)}.
\]
This equivalence unpacks the required sheaves, gluing maps, and
action compatibility. It does not construct any of them.
The local comparison must connect the Ben-Zvi--Chen--Helm--Nadler
construction at \(\ell\ne p\) with the Emerton--Gee target at
\(\ell=p\), including the required coefficient comparison.
The cross-prime construction requires prismatic transition maps,
their triple-overlap cocycle, and continuity of the induced
\(\mathrm{IndCoh}\) construction.
These data remain conjectural in
\Cref{v4-arith:bcglue-k5:thm:prismatic-gluing}.

\paragraph{Global Bernstein-centre compatibility.}
The core condition compares the glued action with the global
Bernstein-centre action under the arithmetic Arinkin--Gaitsgory
identification. That identification is itself a hypothesis.
The reduction of
\Cref{v4-arith:bcglue-k5:cor:LACD-tight-list} also retains
F2.c, the Bernstein lift, local categorical comparison,
the archimedean boundary datum, and global identification from
\Cref{v4-arith:lacd:cor:LACD-residual}.
Supplying the local and cross-prime data leaves the core action
compatibility to be established under those hypotheses.

\subsection{Primary Source Identifiers in App.~A}
\label{v4-arith:w3-comparison-A:appA}

\paragraph{Local geometric and categorical inputs.}
Emerton--Gee arXiv:1908.07185 constructs the local moduli stacks.
Ben-Zvi--Chen--Helm--Nadler arXiv:2010.02321 supplies coherent
Springer theory, and Fargues--Scholze arXiv:2102.13459 supplies
local geometrization and spectral actions.
Dotto--Emerton--Gee arXiv:2603.26887 supplies the fully faithful,
fixed-central-character functor at \(p\geq5\) used in the F1.a
comparison of \Cref{v4-arith:cl:F1-closure} and \Cref{v4-arith:of:F1}.
The full local categorical equivalence requires further input.

\section{Logical Form for F1--F3 Arithmetic}
\label{v4-arith:w3-comparison-A:summary}

The local categorical comparison F1.a requires more than the fully
faithful, fixed-central-character functor at \(p\geq5\).
F1.c defines the Selmer target using the chosen local deformation
conditions. Its derived and Noetherian properties remain conditional
on those local conditions, as in \Cref{v4-arith:thm:F1-reduction}.

The formal descent arguments require F2.c and the
\(\infty\)-categorical Bernstein lift.
Local categorical comparison, the archimedean boundary datum,
and global identification with the automorphic category are
additional inputs to
\Cref{v4-arith:lacd:cor:LACD-residual}.
The three Bernstein-centre requirements are local construction,
cross-prime gluing, and compatibility of the global action.
The core action comparison also assumes arithmetic
Arinkin--Gaitsgory.

The structural splits in
\Cref{v4-arith:sw:thm:LACD-split} and
\Cref{v4-arith:sw:thm:BZN-Ar-split}
organize these comparison problems.
The conditional implications above use their specified construction
data. A converse identifying the original global assertions with
only the abbreviated lists of subconditions requires additional
comparison maps and hypotheses.
The strong-approximation input for \(SL_2\) and the determinant
idele-class quotient for \(GL_2\) remain distinct from these
categorical comparisons.

For the arithmetic trace comparison, the Selmer fibre product and
its tangent complex specify R2. The required formal geometry and
nilpotent-support compatibility remain hypotheses.
R1.formal supplies the extension beyond the finite-type trace theorem.
The loop comparison R3 additionally requires the three completion
and generation assumptions stated above. It leaves the full
Weil/\((\varphi,\Gamma)\) target comparison separate.
R4 includes the wild-prime contributions, with large-prime vanishing
conditional on the arithmetic semisimplicity datum.
Together with F1 and F2, these are the inputs to the conditional
arithmetic Ben-Zvi--Nadler construction.

\section{Open Conditions for F1--F3 Arithmetic}
\label{v4-arith:w3-comparison-A:overall}

The remaining F1--F3 conditions are as follows,
consistent with the running list \Cref{v4-arith:app:C:residual-table-w2}
and the RH completion conjunction \Cref{v4-arith:platonic:def:RH-package}.

\begin{itemize}
\item F1.a (local categorical equivalence
\(\mathrm{D}^{\mathrm{sph\text{-}fin}}([GL_2(\mathbb{Q}_p)])
\simeq
\mathrm{IndCoh}_{\mathrm{Nilp}}(\mathcal{X}^{\mathrm{EG},p}_{GL_2})\)):
\ClaimStatusNeedsVerification\ via Dotto--Emerton--Gee
arXiv:2603.26887 (fully-faithful-only, \(p\geq 5\), fixed central
character). Fargues--Scholze arXiv:2102.13459 supplies the local
geometrisation, spectral action, and Langlands-parameter structure;
the categorical equivalence is the residual.

\item F1 is conditional on F1.a and \(\mathrm{L}_{\mathrm{ACD}}\);
\(\mathrm{L}_{\mathrm{ACD}}\) requires F2.c tightened,
(BCglue.local), (BCglue.glue), and (BC-diamond-glue.core).

\item F3.BZN\({}^{\mathrm{Ar}}\) requires F1, F2, R1.formal.b
(profinite formal-loop continuity),
R1.formal.c (flat-locus pullback),
R4 (wild-prime \(\mathrm{Ext}^1\) refinement, domain beyond
the explicit computation already stated), and the
full-target Weil/\((\varphi,\Gamma)\) comparison.

\item \Cref{v4-arith:platonic:def:RH-package} is the sufficient set
for RH, with \(VB^*\) operator positivity as a forward positivity
residual, \(H\text{-}P^{\mathrm{Ar}}\) as a separate spectral
hypothesis, and the arithmetic Positselski C1--C5 plus
Koszul--Petersson compatibility as the four-input RH form conditions.
\end{itemize}
